% BHU PYQ -- Manually transcribed & error-corrected
% Paper: BCH-212 : Business Economics-I
% Exam: B.Com. (Hons.) Semester III Examination, 2023-24

\documentclass[11pt,a4paper]{article}
\usepackage[a4paper,top=2.5cm,bottom=2.5cm,left=2.8cm,right=2.8cm,headheight=14pt]{geometry}
\usepackage{amsmath,amssymb}
\usepackage[T1]{fontenc}
\usepackage[utf8]{inputenc}
\usepackage{lmodern,microtype,fancyhdr,enumitem,hyperref,lastpage,parskip,booktabs}

\hypersetup{
    colorlinks=true,
    urlcolor=black,
    linkcolor=black,
    pdftitle={BCH-212: Business Economics-I -- BHU B.Com. (Hons.) Sem III 2023-24}
}

\pagestyle{fancy}
\fancyhf{}
\renewcommand{\headrulewidth}{0.4pt}
\renewcommand{\footrulewidth}{0.4pt}
\fancyhead[L]{\small   \textit{Banaras Hindu University}}
\fancyhead[R]{\small   \textit{BCH-212: Business Economics-I}}
\fancyfoot[C]{\small Page  \thepage\ of \pageref{LastPage}}


\newlist{parts}{enumerate}{2}
\setlist[parts,1]{label=(\alph*),leftmargin=*,itemsep=5pt,topsep=3pt}
\setlist[parts,2]{label=(\roman*),leftmargin=*,itemsep=3pt,topsep=2pt}

\newcommand{\pts}[1]{\hfill   \textit{[#1]}}


\begin{document}


\begin{center}
  {\large\bfseries BANARAS HINDU UNIVERSITY}\\[2pt]
  {
\normalsize B.Com. (Hons.) --- Semester III Examination, 2023-24}\\[6pt]
  \rule{0.8\linewidth}{0.5pt}\\[6pt]
  {\large\bfseries Paper No. BCH-212: Business Economics-I}\\[4pt]
  {
\normalsize Time: 3 Hours\hfill Maximum Marks: 70}
\end{center}

\rule{\linewidth}{0.4pt}

\smallskip



\noindent   \textit{Instructions: Answer all questions. All questions carry equal marks (14 marks each).}

\bigskip




\noindent   \textbf{Question 1.}
Discuss the nature and scope of business economics. How is economics different from business economics?\pts{14}

\centerline{   \textbf{OR}}

Describe the relationship of business economics with other branches of knowledge.\pts{14}

\medskip\rule{\linewidth}{0.2pt}




\noindent   \textbf{Question 2.}
Critically analyze the law of equi-marginal utility and mention its importance.\pts{14}

\centerline{   \textbf{OR}}

Describe the meaning and assumptions of Indifference Curves. Explain the state of equilibrium of a consumer with the help of these.\pts{14}

\medskip\rule{\linewidth}{0.2pt}




\noindent   \textbf{Question 3.}
Discuss the law of demand. Why does the demand curve slope down from left to right? Discuss the exceptions to the law of demand.\pts{14}

\centerline{   \textbf{OR}}

Explain the concept of the elasticity of demand. How is the elasticity of demand measured?\pts{14}

\medskip\rule{\linewidth}{0.2pt}




\noindent   \textbf{Question 4.}
Discuss the various methods of Demand forecasting. Mention the steps in Demand Forecasting.\pts{14}

\centerline{   \textbf{OR}}

Write notes on the following:\pts{14}

\begin{parts}
  \item Consumer's surplus
  \item Significance of Demand Forecasting
\end{parts}

\medskip\rule{\linewidth}{0.2pt}




\noindent   \textbf{Question 5.}
Explain the law of variable proportions. Discuss the importance of this law.\pts{14}

\centerline{   \textbf{OR}}

Mention the behaviour of average cost curves during short and long periods.\pts{14}

\vfill
\rule{\linewidth}{0.4pt}

\begin{center}\small
\href{https://bhu-exam.vercel.app/}{Click here to visit our website}\\[4pt]
\href{https://me-aryan.vercel.app/}{Click here to visit our contributor}\\[4pt]
\href{https://quashan-me.vercel.app/}{Click here to visit our contributor}
\end{center}

\end{document}
